\scanpage{173}
% archive sheet number: 172
\diagram{$\begin{array}{ccccccccc|c}
{\unsure{\mathbb G|G}}_m&\mathbb Z/\unsure{\ell|l}\mathbb Z&\unsure{\mu_p|H_p}&(\mathbb G_a&\alpha_p&\unsure{\mathbb V|V})&\mathbb Z/p\mathbb Z&\mathbb Z/\ell\mathbb Z&\mathbb Z&\mathrm{Ab}\\
F\text{ épim}&\left.\begin{matrix}F\\V\end{matrix}\right\}\text{isom}&F=0&F\text{ épim}&F=0&F=0&F\text{ isom}&\left.\begin{matrix}F\\V\end{matrix}\right\}\text{isom}&F\text{ isom}&F\text{ épim}\\
V\text{ isom}&&V\text{ isom}&V=0&V=0&V\text{ \unsure{monom|épimom}}&V=0&&V\text{ monom}&V\text{ n'importe quoi}\\
&&\multicolumn{4}{c}{\text{\begin{tikzpicture}\draw (0,.16)--(0,0)--(4.7,0)--(4.7,.16);\draw (3.1,0)--(3.1,.16);\end{tikzpicture}}} & * &&&
\end{array}$}

\[
\begin{aligned}
F\text{ isom}&:\quad\text{groupe discret séparable}\\
F^n=0&:\quad\text{groupe infinitésimal}
\end{aligned}
\]
\[
F\text{ épim}\quad :\quad \text{groupe \inserted{ind-}\uline{algébrique} séparable}
\]
$\underset{?}{\Longrightarrow}$ (peut avoir un quotient $\mathbb Z$, mais pas de facteur de composition $\mathbb V$)

i.e. pas de partie infinitésimale

\[
\begin{aligned}
V\text{ isom}&:\quad\text{groupe multiplicatif}\\
V^n=0&:\quad\text{groupe \inserted{algébrique} unipotent}
\end{aligned}
\]
\[
V\text{ \deleted{\unsure{épim}}\unsure{monom|épimom}}\quad :\quad\deleted{\text{\unsure{groupe inf.}}}\quad\text{\inserted{algébrique}}
\]
$\overset{?}{\Longrightarrow}$ pas de partie unipotente en facteur de composition

[Condition de finitude est \unsure{superflue} \unsure{pour|si on} une extension de groupe formel par groupe algébrique]

[gr. infinitésimal $\ne$ groupe radiciel

$\mathbb V$ n'est pas unipotent ? bien qu'il soit une extension unipotente

\[
\left.\begin{aligned}
F\text{ monom}&\Longrightarrow F\text{ isom}\\
V\text{ épim}&\Longrightarrow V\text{ isom}
\end{aligned}\right\}?
\]

Il y a 8 types fondamentaux de groupes affines (satisfaisant la condition de finitude)

Le passage aux limites inductives resp. projectives introduit d'autres types \inserted{unipotents} (les types de Dieudonné). Je pense que les tableaux ci-dessous restent valables essentiellement.

\margin{\unsure{Ça|Ca} pourra faire \unsure{apparaître|apparaitre} d'autres types de dualité, liés à la dualité des var. abéliennes \unsure{et aux} $\unsure{p^\infty|p}$, à la théorie des dualités des \unsure{modules de Cartier|modules de Cartan} \unsure{par}.}
\diagram{\begin{tikzcd}[ampersand replacement=\&,column sep=5em]
{}\arrow[r,leftrightarrow]\&{}
\end{tikzcd}}
